\documentclass[../main.tex]{subfiles}

\begin{document}

\chapter{Alineamiento}

\section{Modo de alineamiento Flat}

\begin{figure}[htbp]
    \centering
    \includegraphics[width=1\textwidth]{Alignment/FlatAlignment.png}
    \caption{Pendiente.}
    %\label{fig:frontDA1000}
\end{figure}


\begin{longtblr}[
    caption={Parámetros de configuración},
    label={tab:parameters},
]{
    width=\textwidth,
    colspec={|Q[c]|X|Q[c]|},
    hlines,
}

\SetRow{bg=gray9}
\textbf{Parámetro} &
\textbf{Descripción} &
\textbf{Tipo .NET} \\

%===========================================================
% Alignment
%===========================================================
\SetRow{bg=SkyBlue}
\SetCell[c=3]{c}
Alignment \\

mainStageNumberX &
Número del eje principal utilizado para el movimiento en X. &
UInt16 \\

mainStageNumberY &
Número del eje principal utilizado para el movimiento en Y. &
UInt16 \\

subStageNumberXY & 
Número del sub-eje utilizado para el movimiento en XY. & 
UInt16 \\

subAngleX & 
Ángulo del sub-eje para la alineación en X. & 
Double \\

subAngleY & 
Ángulo del sub-eje para la alineación en Y. & 
Double \\

%===========================================================
% Sensor
%===========================================================
\SetRow{bg=SweetSalmon}
\SetCell[c=3]{c}
Sensor \\

pmCh & Canal del medidor de potencia utilizado durante la alineación. & UInt16 \\

analogCh & Canal de entrada analógica utilizado para la medición. & UInt16 \\

pmAutoRangeUpOn & Habilita el incremento automático del rango del medidor de potencia. & Boolean \\

pmInitRangeSettingOn & Habilita el uso de un rango inicial para el medidor de potencia. & Boolean \\

pmInitRange & Valor del rango inicial del medidor de potencia. & Int16 \\

%===========================================================
% Threshold
%===========================================================
\SetRow{bg=Periwinkle}
\SetCell[c=3]{c}
Threshold \\

fieldSearchThreshold & Umbral utilizado durante la búsqueda de campo. & Double \\

peakSearchThreshold & Umbral utilizado durante la búsqueda del pico de señal. & Double \\

%===========================================================
% Search
%===========================================================
\SetRow{bg=MutedTeal}
\SetCell[c=3]{c}
Search \\

searchRangeX & Rango de búsqueda en la dirección X. & Double \\

searchRangeY & Rango de búsqueda en la dirección Y. & Double \\

fieldSearchPitchX & Paso de exploración en X durante la búsqueda de campo. & Double \\

fieldSearchPitchY & Paso de exploración en Y durante la búsqueda de campo. & Double \\

fieldSearchFirstPitchX & Primer paso de exploración en X durante la búsqueda de campo. & Double \\

fieldSearchSpeedX & Velocidad de exploración en X durante la búsqueda de campo. & Double \\

fieldSearchSpeedY & Velocidad de exploración en Y durante la búsqueda de campo. & Double \\

peakSearchSpeedX & Velocidad de búsqueda del pico en X. & Double \\

peakSearchSpeedY & Velocidad de búsqueda del pico en Y. & Double \\

smoothingRangeX & Ventana de suavizado aplicada a los datos en X. & UInt32 \\

smoothingRangeY & Ventana de suavizado aplicada a los datos en Y. & UInt32 \\

centroidThresholdX & Umbral para el cálculo del centroide en X. & Double \\

centroidThresholdY & Umbral para el cálculo del centroide en Y. & Double \\

%===========================================================
% Judge
%===========================================================
\SetRow{bg=ApricotCream}
\SetCell[c=3]{c}
Judge \\

convergentRangeX & Rango de convergencia permitido en X. & Double \\

convergentRangeY & Rango de convergencia permitido en Y. & Double \\

comparisonCount & Número de comparaciones realizadas para verificar la convergencia. & UInt16 \\

maxRepeatCount & Número máximo de repeticiones permitidas durante el proceso de alineación. & UInt16 \\

\end{longtblr}










% \begin{longtblr}[
%     caption={Parámetros de configuración},
%     label={tab:parameters},
% ]{
%     width=\textwidth,
%     colspec={|Q[c]|X|Q[c]|},
%     hlines,
% }

% \SetRow{bg=gray9}
% \textbf{Parámetro} &
% \textbf{Descripción} &
% \textbf{Tipo .NET} \\

% Parámetro &
% \SetCell[c=2]{c}
% Texto de ejemplo &
% \\

% mainStageNumberX &
% Número del eje principal utilizado para el movimiento en X. &
% UInt16 \\

% mainStageNumberY &
% Número del eje principal utilizado para el movimiento en Y. &
% UInt16 \\

% maxRepeatCount &
% Número máximo de repeticiones permitidas durante el proceso de alineación. &
% UInt16 \\

% mainStageNumberX & Número del eje principal utilizado para el movimiento en X. & UInt16 \\

% mainStageNumberY & Número del eje principal utilizado para el movimiento en Y. & UInt16 \\

% subStageNumberXY & Número del sub-eje utilizado para el movimiento en XY. & UInt16 \\

% subAngleX & Ángulo del sub-eje para la alineación en X. & Double \\

% subAngleY & Ángulo del sub-eje para la alineación en Y. & Double \\

% pmCh & Canal del medidor de potencia utilizado durante la alineación. & UInt16 \\

% analogCh & Canal de entrada analógica utilizado para la medición. & UInt16 \\

% pmAutoRangeUpOn & Habilita el incremento automático del rango del medidor de potencia. & Boolean \\

% pmInitRangeSettingOn & Habilita el uso de un rango inicial para el medidor de potencia. & Boolean \\

% pmInitRange & Valor del rango inicial del medidor de potencia. & Int16 \\

% fieldSearchThreshold & Umbral utilizado durante la búsqueda de campo. & Double \\

% peakSearchThreshold & Umbral utilizado durante la búsqueda del pico de señal. & Double \\

% searchRangeX & Rango de búsqueda en la dirección X. & Double \\

% searchRangeY & Rango de búsqueda en la dirección Y. & Double \\

% fieldSearchPitchX & Paso de exploración en X durante la búsqueda de campo. & Double \\

% fieldSearchPitchY & Paso de exploración en Y durante la búsqueda de campo. & Double \\

% fieldSearchFirstPitchX & Primer paso de exploración en X durante la búsqueda de campo. & Double \\

% fieldSearchSpeedX & Velocidad de exploración en X durante la búsqueda de campo. & Double \\

% fieldSearchSpeedY & Velocidad de exploración en Y durante la búsqueda de campo. & Double \\

% peakSearchSpeedX & Velocidad de búsqueda del pico en X. & Double \\

% peakSearchSpeedY & Velocidad de búsqueda del pico en Y. & Double \\

% smoothingRangeX & Ventana de suavizado aplicada a los datos en X. & UInt32 \\

% smoothingRangeY & Ventana de suavizado aplicada a los datos en Y. & UInt32 \\

% centroidThresholdX & Umbral para el cálculo del centroide en X. & Double \\

% centroidThresholdY & Umbral para el cálculo del centroide en Y. & Double \\

% convergentRangeX & Rango de convergencia permitido en X. & Double \\

% convergentRangeY & Rango de convergencia permitido en Y. & Double \\

% comparisonCount & Número de comparaciones realizadas para verificar la convergencia. & UInt16 \\

% maxRepeatCount & Número máximo de repeticiones permitidas durante el proceso de alineación. & UInt16 \\


% \end{longtblr}

% \begin{longtblr}[H]
% \centering
% \resizebox{\textwidth}{!}{%
% \begin{tabular}{|c|p{8cm}|c|}
% \hline
%  Parametro & Descripción & Tipo .NET \\
% \hline
% mainStageNumberX & Número del eje principal utilizado para el movimiento en X. & UInt16 \\
% \hline
% mainStageNumberY & Número del eje principal utilizado para el movimiento en Y. & UInt16 \\
% \hline
% subStageNumberXY & Número del sub-eje utilizado para el movimiento en XY. & UInt16 \\
% \hline
% subAngleX & Ángulo del sub-eje para la alineación en X. & Double \\
% \hline
% subAngleY & Ángulo del sub-eje para la alineación en Y. & Double \\
% \hline
% pmCh & Canal del medidor de potencia utilizado durante la alineación. & UInt16 \\
% \hline
% analogCh & Canal de entrada analógica utilizado para la medición. & UInt16 \\
% \hline
% pmAutoRangeUpOn & Habilita el incremento automático del rango del medidor de potencia. & Boolean \\
% \hline
% pmInitRangeSettingOn & Habilita el uso de un rango inicial para el medidor de potencia. & Boolean \\
% \hline
% pmInitRange & Valor del rango inicial del medidor de potencia. & Int16 \\
% \hline
% fieldSearchThreshold & Umbral utilizado durante la búsqueda de campo. & Double \\
% \hline
% peakSearchThreshold & Umbral utilizado durante la búsqueda del pico de señal. & Double \\
% \hline
% searchRangeX & Rango de búsqueda en la dirección X. & Double \\
% \hline
% searchRangeY & Rango de búsqueda en la dirección Y. & Double \\
% \hline
% fieldSearchPitchX & Paso de exploración en X durante la búsqueda de campo. & Double \\
% \hline
% fieldSearchPitchY & Paso de exploración en Y durante la búsqueda de campo. & Double \\
% \hline
% fieldSearchFirstPitchX & Primer paso de exploración en X durante la búsqueda de campo. & Double \\
% \hline
% fieldSearchSpeedX & Velocidad de exploración en X durante la búsqueda de campo. & Double \\
% \hline
% fieldSearchSpeedY & Velocidad de exploración en Y durante la búsqueda de campo. & Double \\
% \hline
% peakSearchSpeedX & Velocidad de búsqueda del pico en X. & Double \\
% \hline
% peakSearchSpeedY & Velocidad de búsqueda del pico en Y. & Double \\
% \hline
% smoothingRangeX & Ventana de suavizado aplicada a los datos en X. & UInt32 \\
% \hline
% smoothingRangeY & Ventana de suavizado aplicada a los datos en Y. & UInt32 \\
% \hline
% centroidThresholdX & Umbral para el cálculo del centroide en X. & Double \\
% \hline
% centroidThresholdY & Umbral para el cálculo del centroide en Y. & Double \\
% \hline
% convergentRangeX & Rango de convergencia permitido en X. & Double \\
% \hline
% convergentRangeY & Rango de convergencia permitido en Y. & Double \\
% \hline
% comparisonCount & Número de comparaciones realizadas para verificar la convergencia. & UInt16 \\
% \hline
% maxRepeatCount & Número máximo de repeticiones permitidas durante el proceso de alineación. & UInt16 \\
% \hline
% \end{tabular}
% } % End Resizebox
% \end{longtblr}





% \begin{table}[H]
% \centering
% \resizebox{\textwidth}{!}{%
% \begin{tabular}{|c|p{8cm}|c|}
% \hline
%  Parametro & Descripción & Tipo .NET \\
% \hline
% mainStageNumberX & Número del eje principal utilizado para el movimiento en X. & UInt16 \\
% \hline
% mainStageNumberY & Número del eje principal utilizado para el movimiento en Y. & UInt16 \\
% \hline
% subStageNumberXY & Número del sub-eje utilizado para el movimiento en XY. & UInt16 \\
% \hline
% subAngleX & Ángulo del sub-eje para la alineación en X. & Double \\
% \hline
% subAngleY & Ángulo del sub-eje para la alineación en Y. & Double \\
% \hline
% pmCh & Canal del medidor de potencia utilizado durante la alineación. & UInt16 \\
% \hline
% analogCh & Canal de entrada analógica utilizado para la medición. & UInt16 \\
% \hline
% pmAutoRangeUpOn & Habilita el incremento automático del rango del medidor de potencia. & Boolean \\
% \hline
% pmInitRangeSettingOn & Habilita el uso de un rango inicial para el medidor de potencia. & Boolean \\
% \hline
% pmInitRange & Valor del rango inicial del medidor de potencia. & Int16 \\
% \hline
% fieldSearchThreshold & Umbral utilizado durante la búsqueda de campo. & Double \\
% \hline
% peakSearchThreshold & Umbral utilizado durante la búsqueda del pico de señal. & Double \\
% \hline
% searchRangeX & Rango de búsqueda en la dirección X. & Double \\
% \hline
% searchRangeY & Rango de búsqueda en la dirección Y. & Double \\
% \hline
% fieldSearchPitchX & Paso de exploración en X durante la búsqueda de campo. & Double \\
% \hline
% fieldSearchPitchY & Paso de exploración en Y durante la búsqueda de campo. & Double \\
% \hline
% fieldSearchFirstPitchX & Primer paso de exploración en X durante la búsqueda de campo. & Double \\
% \hline
% fieldSearchSpeedX & Velocidad de exploración en X durante la búsqueda de campo. & Double \\
% \hline
% fieldSearchSpeedY & Velocidad de exploración en Y durante la búsqueda de campo. & Double \\
% \hline
% peakSearchSpeedX & Velocidad de búsqueda del pico en X. & Double \\
% \hline
% peakSearchSpeedY & Velocidad de búsqueda del pico en Y. & Double \\
% \hline
% smoothingRangeX & Ventana de suavizado aplicada a los datos en X. & UInt32 \\
% \hline
% smoothingRangeY & Ventana de suavizado aplicada a los datos en Y. & UInt32 \\
% \hline
% centroidThresholdX & Umbral para el cálculo del centroide en X. & Double \\
% \hline
% centroidThresholdY & Umbral para el cálculo del centroide en Y. & Double \\
% \hline
% convergentRangeX & Rango de convergencia permitido en X. & Double \\
% \hline
% convergentRangeY & Rango de convergencia permitido en Y. & Double \\
% \hline
% comparisonCount & Número de comparaciones realizadas para verificar la convergencia. & UInt16 \\
% \hline
% maxRepeatCount & Número máximo de repeticiones permitidas durante el proceso de alineación. & UInt16 \\
% \hline
% \end{tabular}
% } % End Resizebox
% \end{table}


% % Hola

% \begin{table}[H]
% \centering
% \resizebox{\textwidth}{!}{%
% \begin{tabularx}{\textwidth}{|c|X|c|}
% \hline
% \textbf{Parámetro} & \textbf{Descripción} & \textbf{Tipo .NET} \\
% \hline
% mainStageNumberX & Número del eje principal utilizado para el movimiento en X. & UInt16 \\
% \hline
% mainStageNumberY & Número del eje principal utilizado para el movimiento en Y. & UInt16 \\
% \hline
% subStageNumberXY & Número del sub-eje utilizado para el movimiento en XY. & UInt16 \\
% \hline
% subAngleX & Ángulo del sub-eje para la alineación en X. & Double \\
% \hline
% subAngleY & Ángulo del sub-eje para la alineación en Y. & Double \\
% \hline
% pmCh & Canal del medidor de potencia utilizado durante la alineación. & UInt16 \\
% \hline
% analogCh & Canal de entrada analógica utilizado para la medición. & UInt16 \\
% \hline
% pmAutoRangeUpOn & Habilita el incremento automático del rango del medidor de potencia. & Boolean \\
% \hline
% pmInitRangeSettingOn & Habilita el uso de un rango inicial para el medidor de potencia. & Boolean \\
% \hline
% pmInitRange & Valor del rango inicial del medidor de potencia. & Int16 \\
% \hline
% fieldSearchThreshold & Umbral utilizado durante la búsqueda de campo. & Double \\
% \hline
% peakSearchThreshold & Umbral utilizado durante la búsqueda del pico de señal. & Double \\
% \hline
% searchRangeX & Rango de búsqueda en la dirección X. & Double \\
% \hline
% searchRangeY & Rango de búsqueda en la dirección Y. & Double \\
% \hline
% fieldSearchPitchX & Paso de exploración en X durante la búsqueda de campo. & Double \\
% \hline
% fieldSearchPitchY & Paso de exploración en Y durante la búsqueda de campo. & Double \\
% \hline
% fieldSearchFirstPitchX & Primer paso de exploración en X durante la búsqueda de campo. & Double \\
% \hline
% fieldSearchSpeedX & Velocidad de exploración en X durante la búsqueda de campo. & Double \\
% \hline
% fieldSearchSpeedY & Velocidad de exploración en Y durante la búsqueda de campo. & Double \\
% \hline
% peakSearchSpeedX & Velocidad de búsqueda del pico en X. & Double \\
% \hline
% peakSearchSpeedY & Velocidad de búsqueda del pico en Y. & Double \\
% \hline
% smoothingRangeX & Ventana de suavizado aplicada a los datos en X. & UInt32 \\
% \hline
% smoothingRangeY & Ventana de suavizado aplicada a los datos en Y. & UInt32 \\
% \hline
% centroidThresholdX & Umbral para el cálculo del centroide en X. & Double \\
% \hline
% centroidThresholdY & Umbral para el cálculo del centroide en Y. & Double \\
% \hline
% convergentRangeX & Rango de convergencia permitido en X. & Double \\
% \hline
% convergentRangeY & Rango de convergencia permitido en Y. & Double \\
% \hline
% comparisonCount & Número de comparaciones realizadas para verificar la convergencia. & UInt16 \\
% \hline
% maxRepeatCount & Número máximo de repeticiones permitidas durante el proceso de alineación. & UInt16 \\
% \hline
% \end{tabularx}
% } % End Resizebox
% \end{table}



% \begin{table}[H]
% \centering
% \begin{tabularx}{\textwidth}{|c|X|c|}
% \hline
%  Parametro & Descripción & Tipo .NET \\
% \hline
% %--------------------------------------
% mainStageNumberX &
% Pendiente &
% UInt16 \\
% \hline
% \end{tabularx}
% \end{table}



\end{document}



